\chapter{مبرهنة باربان--دافنبورت--هالبرستام ومتوسط التباين التربيعي}

\section{نطاق الفصل وحالته}

\noindent\textbf{حالة الفصل:} \textenglish{AUTHORED-DRAFT}. أُنشئ هذا المتن بعد
إغلاق بوابة ما قبل التأليف بحكم مستقل
\textenglish{APPROVED-WITH-NONBLOCKING-CORRECTIONS}، ثم صدور قرار المالك
\textenglish{PASS-FOR-AUTHORING}. لا يزال الفصل خاضعًا لتدقيق ما بعد التأليف
ولفحوص الجودة وبناء ملف المعاينة.

يبحث الفصل متوسط مربع خطأ توزيع الأعداد الأولية في المتتاليات الحسابية. إذا
كتبنا
\[
\psi(x;q,a)=\sum_{\substack{n\le x\\n\equiv a\pmod q}}\Lambda(n),
\]
فإن الكمية المركزية هي
\[
V_\psi(x,Q)=
\sum_{q\le Q}
\sum_{\substack{a\bmod q\\(a,q)=1}}
\left|\psi(x;q,a)-\frac{x}{\varphi(q)}\right|^2.
\]
سنثبت أنه لكل ثابت \(A>0\)، وبانتظام عندما
\[
x\ge3,
\qquad
\frac{x}{(\log x)^A}\le Q\le x,
\]
لدينا
\[
V_\psi(x,Q)\ll_A xQ\log x.
\]
الثابت الضمني يعتمد على \(A\)، وهو غير فعال في المسار الحالي بسبب مدخل
سيغل--فالفش للموصلات الصغيرة.

لا يثبت هذا الفصل الصيغة التقاربية لمونتغمري--هولي، ولا مبرهنة باربان العامة
في كل مجالات \(Q\)، ولا أي نتيجة مشروطة بفرضية ريمان المعممة. كما لا يستعمل
مبرهنة بومبييري--فينوغرادوف في البرهان، رغم قرب الموضوعين مفهوميًا.

\subsection{نواتج التعلم}

بعد إتمام الفصل ينبغي أن يستطيع القارئ:
\begin{enumerate}
  \item تحويل تباين الفئات المختزلة إلى متوسط تربيعي على شخصيات ديريشليه.
  \item رد الشخصيات المستحثة إلى موصلاتها مع ضبط التصحيحات المحلية.
  \item فهم وزن الموصل
  \[
  W_Q(r)=\sum_{m\le Q/r}\frac1{\varphi(rm)}.
  \]
  \item استعمال الغربال الكبير للشخصيات البدائية على كتل ديادية.
  \item تفسير اختيار حد الفصل
  \[
  Y=\max\!\left(2,\frac{x}{Q}\log(2Q)\right).
  \]
  \item فصل الموصلات الصغيرة التي تعالج بسيغل--فالفش عن الموصلات الكبيرة
  التي يعالجها الغربال الكبير مباشرة.
  \item تتبع الشخصية الرئيسية والموصل \(1\) وعدم فعالية الثابت النهائي.
  \item التمييز بين الحد العلوي الكلاسيكي المثبت هنا والصيغ الأقوى المؤجلة.
\end{enumerate}

\section{التوسيط وتحويل التباين إلى الشخصيات}

لكل شخصية ديريشليه \(\chi\bmod q\) نضع
\[
\Psi(x,\chi)=\sum_{n\le x}\Lambda(n)\chi(n)
\]
ونعرّف المجموع المتمركز
\[
\Psi^\circ(x,\chi)
=
\Psi(x,\chi)-x\,\mathbf 1_{\chi=\chi_0^{(q)}}.
\]
إذن الشخصية غير الرئيسية لا تتغير، بينما
\[
\Psi^\circ(x,\chi_0^{(q)})
=
\sum_{\substack{n\le x\\(n,q)=1}}\Lambda(n)-x.
\]

\begin{lemma}[هوية التباين بالشخصيات]
\label{lem:chapter14-character-variance}
\resultid{ANT-LEM-14-01}
\provedhere
لكل \(q\ge1\):
\[
\sum_{\substack{a\bmod q\\(a,q)=1}}
\left|\psi(x;q,a)-\frac{x}{\varphi(q)}\right|^2
=
\frac1{\varphi(q)}
\sum_{\chi\bmod q}|\Psi^\circ(x,\chi)|^2.
\tag{14.1}
\]
\end{lemma}

\begin{proof}
بتعامد شخصيات المجموعة المختزلة:
\[
\psi(x;q,a)
=
\frac1{\varphi(q)}
\sum_{\chi\bmod q}\overline{\chi(a)}\Psi(x,\chi).
\]
وبما أن الشخصية الرئيسية تساوي \(1\) على الفئات المختزلة، فإن طرح
\(x/\varphi(q)\) يعادل استبدال \(\Psi\) بـ\(\Psi^\circ\). وبعد تربيع القيمة
المطلقة والجمع على \(a\)، تلغي علاقة التعامد الحدود المختلطة وتعطي العامل
\(1/\varphi(q)\) بالضبط.
\end{proof}

بالتالي يكفي ضبط
\[
\mathcal S(x,Q)
=
\sum_{q\le Q}\frac1{\varphi(q)}
\sum_{\chi\bmod q}|\Psi^\circ(x,\chi)|^2,
\]
إذ إن
\[
V_\psi(x,Q)=\mathcal S(x,Q).
\tag{14.2}
\]

\section{رد الشخصيات إلى الموصلات}

لكل شخصية \(\chi\bmod q\) موصل وحيد \(r\mid q\) وشخصية بدائية وحيدة
\(\chi^*\bmod r\) تستحثها. نكتب \(q=rm\)، ونعرّف
\[
W_Q(r)=\sum_{m\le Q/r}\frac1{\varphi(rm)}.
\]

\begin{lemma}[حد وزن الموصل]
\label{lem:chapter14-conductor-weight}
\resultid{ANT-LEM-14-02}
\provedhere
بانتظام لكل \(1\le r\le Q\):
\[
\frac1{\varphi(r)}
\le W_Q(r)
\ll
\frac{\log(2Q/r)}{\varphi(r)}.
\tag{14.3}
\]
\end{lemma}

\begin{proof}
لدينا
\[
\varphi(rm)\ge\varphi(r)\varphi(m),
\]
ومن ثم
\[
W_Q(r)
\le
\frac1{\varphi(r)}
\sum_{m\le Q/r}\frac1{\varphi(m)}.
\]
وباستعمال الهوية
\[
\frac{n}{\varphi(n)}
=
\sum_{d\mid n}\frac{\mu^2(d)}{\varphi(d)}
\]
والجمع التوافقي، نحصل على
\[
\sum_{m\le M}\frac1{\varphi(m)}\ll\log(2M).
\]
أما الحد السفلي فيأتي من الحد \(m=1\).
\end{proof}

إذا كانت \(\chi\bmod q\) مستحثة من \(\chi^*\bmod r\)، فإن الفرق بين
المجموعين مدعوم على القوى الأولية التي تقسم \(q\) ولا تقسم \(r\):
\[
\Psi^\circ(x,\chi)
=
\Psi^\circ(x,\chi^*)-C(x;q,r,\chi^*),
\]
حيث
\[
C(x;q,r,\chi^*)
=
\sum_{\substack{p^k\le x\\p\mid q,\ p\nmid r}}
\chi^*(p^k)\log p,
\qquad
|C(x;q,r,\chi^*)|\le\omega(m)\log x.
\tag{14.4}
\]

\begin{lemma}[كلفة تصحيحات الاستحثاث]
\label{lem:chapter14-local-corrections}
\resultid{ANT-LEM-14-03}
\provedhere
لدينا
\[
\sum_{q\le Q}\frac1{\varphi(q)}
\sum_{\chi\bmod q}|C(x;q,r_\chi,\chi^*)|^2
\ll Q(\log x)^2.
\tag{14.5}
\]
وبالتالي، إذا عرّفنا
\[
\mathcal P(x,Q)
=
\sum_{r\le Q}W_Q(r)
\sum_{\chi^*\bmod r}^{*}|\Psi^\circ(x,\chi^*)|^2,
\]
فإن
\[
\mathcal S(x,Q)
\le2\mathcal P(x,Q)+O\!\left(Q(\log x)^2\right).
\tag{14.6}
\]
\end{lemma}

\begin{proof}
بإعادة الفهرسة \(q=rm\)، وعدد الشخصيات البدائية
\(\varphi^*(r)\le\varphi(r)\)، نحصل على
\[
\mathcal E(x,Q)
\le
(\log x)^2
\sum_{m\le Q}\omega(m)^2
\sum_{r\le Q/m}\frac{\varphi^*(r)}{\varphi(rm)}.
\]
ومن \(\varphi(rm)\ge\varphi(r)\varphi(m)\):
\[
\mathcal E(x,Q)
\le
Q(\log x)^2
\sum_{m\le Q}\frac{\omega(m)^2}{m\varphi(m)}.
\]
وتتقارب المتسلسلة الأخيرة. ولتوضيح ذلك، ضع
\(g(n)=1/(n\varphi(n))\). لكل أولي \(p\):
\[
G_p=\sum_{k\ge1}g(p^k)
=
\frac{p}{(p-1)(p^2-1)}\ll p^{-2}.
\]
وباستعمال
\[
\omega(n)^2
=
\sum_{p\mid n}1
+2\sum_{p<\ell\atop p\ell\mid n}1,
\]
ثم تونيللي والتضاعف، تسيطر المتسلسلة المطلوبة بواسطة ثابت مضروب في
\[
\sum_pG_p+2\sum_{p<\ell}G_pG_\ell<\infty.
\]
أما (14.6) فتنتج من \(|A-B|^2\le2|A|^2+2|B|^2\).
\end{proof}

\section{الغربال الكبير الموزون}

سنستعمل حزمة الغربال الكبير المثبتة مرجعيًا في الفصل الثالث عشر، ولا سيما
المبرهنة \ref{thm:chapter13-large-sieve-package}. الصيغة الأساسية هي
\[
\sum_{r\le R}\frac r{\varphi(r)}
\sum_{\chi\bmod r}^{*}
\left|\sum_{M<n\le M+N}c_n\chi(n)\right|^2
\ll
(N+R^2)\sum|c_n|^2.
\]

\begin{lemma}[الغربال الكبير بوزن الموصل]
\label{lem:chapter14-weighted-large-sieve}
\resultid{ANT-LEM-14-04}
\citedresult{\cite{Bombieri1965LargeSieve,MontgomeryVaughan2026}}
على كتلة \(R<r\le2R\)، ولكل متتالية مركبة \(c_n\) مدعومة في فترة طولها
\(N\):
\[
\sum_{R<r\le2R}W_Q(r)
\sum_{\chi\bmod r}^{*}
\left|\sum_n c_n\chi(n)\right|^2
\ll
\left(\frac NR+R\right)
\log\frac{2Q}{R}
\sum_n|c_n|^2.
\tag{14.7}
\]
\end{lemma}

\begin{proof}
على الكتلة \(R<r\le2R\)، يعطينا (14.3)
\[
W_Q(r)
\ll
\frac{\log(2Q/R)}{\varphi(r)}
\le
\frac{\log(2Q/R)}{R}\frac r{\varphi(r)}.
\]
ثم نطبق الغربال الكبير حتى \(2R\)، فنحصل على
\[
\frac{\log(2Q/R)}{R}(N+4R^2)\sum|c_n|^2,
\]
وهو (14.7).
\end{proof}

\section{فصل الموصلات}

ضع
\[
Y=\max\!\left(2,\frac{x}{Q}\log(2Q)\right).
\tag{14.8}
\]
في المجال
\[
\frac{x}{(\log x)^A}\le Q\le x
\]
لدينا، بعد تكبير حد يعتمد على \(A\):
\[
2\le Y\le Q,
\qquad
Y\ll_A(\log x)^{A+1}.
\tag{14.9}
\]
نفصل
\[
\mathcal P(x,Q)=\mathcal P_{\le Y}(x,Q)+\mathcal P_{>Y}(x,Q).
\]

\subsection{الموصلات الصغيرة}

من مبرهنة سيغل--فالفش في الفصل الثاني عشر، لكل قوة لوغاريتمية ثابتة، وبانتظام
للشخصيات البدائية ذات \(r\le Y\):
\[
|\Psi^\circ(x,\chi^*)|
\ll_A
x\exp(-c_A\sqrt{\log x}).
\]
وهذا يشمل الموصل \(1\)، حيث
\[
\Psi^\circ(x,\chi_1^*)=\psi(x)-x.
\]
وبما أن عدد الشخصيات البدائية modulo \(r\) لا يتجاوز \(\varphi(r)\)، فإن
\[
\mathcal P_{\le Y}(x,Q)
\ll_A
Y\log(2Q)x^2e^{-2c_A\sqrt{\log x}}
=o_A(xQ\log x).
\tag{14.10}
\]
هنا تحديدًا تدخل عدم فعالية الثابت النهائي.

\subsection{الموصلات الكبيرة}

نقسم المجال \((Y,Q]\) إلى الكتل غير المتراكبة
\[
I_j=(R_j,\min(2R_j,Q)],
\qquad R_j=2^jY,
\]
لكل \(j\ge0\) مع \(R_j<Q\).

في المتوسط البدائي لا توجد شخصية رئيسية بدائية ذات موصل \(r>1\): الشخصية
الرئيسية modulo \(q>1\) مستحثة من الموصل \(1\). لذلك، على طبقة
\(r>Y\ge2\):
\[
\Psi^\circ(x,\chi^*)
=
\sum_{n\le x}\Lambda(n)\chi^*(n).
\tag{14.11}
\]
وهذا يجيز تطبيق لمّة \ref{lem:chapter14-weighted-large-sieve} مباشرة على
\(c_n=\Lambda(n)\mathbf1_{n\le x}\).

نستعمل الحد
\[
\sum_{n\le x}\Lambda(n)^2
\le
(\log x)\sum_{n\le x}\Lambda(n)
\ll x\log x,
\tag{14.12}
\]
حيث الحد الأخير هو حد تشيبيشيف من الفصل التاسع. ومن ثم مساهمة الكتلة ذات
المقياس \(R\) لا تتجاوز
\[
\ll
x\log x
\left(\frac{x}{R}+R\right)
\log\frac{2Q}{R}.
\]
أما الجمع الديادي فيعطي
\[
\sum_jR_j\log\frac{2Q}{R_j}\ll Q
\]
و
\[
\sum_j\frac{x}{R_j}\log\frac{2Q}{R_j}
\ll
\frac{x}{Y}\log\frac{2Q}{Y}
\le Q.
\]
لذلك
\[
\mathcal P_{>Y}(x,Q)\ll xQ\log x.
\tag{14.13}
\]
وبجمع (14.10) و(14.13):
\[
\mathcal P(x,Q)\ll_A xQ\log x.
\tag{14.14}
\]

\section{مبرهنة باربان--دافنبورت--هالبرستام في المجال الكلاسيكي}

\begin{theorem}[الحد العلوي الكلاسيكي لباربان--دافنبورت--هالبرستام]
\label{thm:chapter14-bdh}
\resultid{ANT-THM-14-01}
\provedhere
لكل ثابت \(A>0\)، يوجد ثابت \(C_A>0\) بحيث لكل \(x\ge3\) ولكل \(Q\) يحقق
\[
\frac{x}{(\log x)^A}\le Q\le x,
\]
لدينا
\[
\boxed{
\sum_{q\le Q}
\sum_{\substack{a\bmod q\\(a,q)=1}}
\left|\psi(x;q,a)-\frac{x}{\varphi(q)}\right|^2
\le
C_AxQ\log x
}.
\tag{14.15}
\]
الثابت \(C_A\) غير فعال في المسار الحالي.
\end{theorem}

\begin{proof}
من لمّة \ref{lem:chapter14-local-corrections} و(14.14):
\[
\mathcal S(x,Q)
\ll_A
xQ\log x+Q(\log x)^2.
\]
ولأن \(x\ge3\)، يمتص الحد الأول الثاني بعد تعديل الثابت. ثم نستعمل الهوية
(14.2).

البرهان السابق يفترض ضمنيًا أن \(x\) أكبر من حد \(x_0(A)\) يكفل (14.9).
أما المجال المحدود \(3\le x<x_0(A)\)، فيمتص بتكبير الثابت \(C_A\)، لأن
النسبة بين الطرف الأيسر و\(xQ\log x\) محدودة على هذا المجال المحدود.
\end{proof}

\begin{remark}[طبيعة حد الفصل]
الاختيار (14.8) ليس تجميليًا؛ فهو يجعل مساهمة الحد \(x/R\) في الغربال الكبير
بعد أول كتلة من رتبة \(Q\). وفي الوقت نفسه يبقي \(Y\) داخل قوة ثابتة من
\(\log x\)، وهي بالضبط المنطقة التي يغطيها سيغل--فالفش.
\end{remark}

\begin{remark}[المقارنة مع بومبييري--فينوغرادوف]
مبرهنة بومبييري--فينوغرادوف تضبط متوسطًا من النمط \(L^1\) لأكبر خطأ في كل
ترديد حتى مستوى يقارب \(x^{1/2}\). أما (14.15) فتضبط متوسطًا تربيعيًا على
جميع الفئات المختزلة، في مجال \(Q\) القريب من \(x\). لا تتبع إحدى الصيغتين
من الأخرى مباشرة بالتطبيع المستعمل هنا، ولا يستعمل هذا الفصل نتيجة الفصل
الثالث عشر في البرهان.
\end{remark}

\section{المسار المؤجل وعائق التصادمات الضربية}

من الطبيعي محاولة تفكيك \(\Lambda\) بهوية فوغان ثم تطبيق الغربال الكبير على
مستطيلات ثنائية. إذا جُمعت معاملات المستطيل في
\[
c_n=\sum_{mk=n}\alpha_m\beta_k,
\]
فإن القطر الكامل هو \(\sum_n|c_n|^2\)، وليس فقط حاصل معياري
\(\alpha\) و\(\beta\). وبكوشي:
\[
\sum_n|c_n|^2
\le
D_{M,K}\|\alpha\|_2^2\|\beta\|_2^2,
\]
حيث \(D_{M,K}\) أكبر عدد لتمثيلات \(n=mk\) في المستطيل. الحد العام المتاح
من دالة القواسم هو \(x^{o(1)}\)، ولا يعطي خسارة لوغاريتمية ثابتة تكفي للرتبة
الدقيقة المستهدفة.

لهذا لم يخف المشروع العائق، بل عزل مسار فوغان بوصفه طبقة بحثية مؤجلة، واعتمد
بدلًا منه فصل الموصلات في البرهان الكلاسيكي. لا يثبت هذا أن مسار فوغان
مستحيل؛ بل يثبت أن التجميع المباشر قبل استغلال البنية الثنائية غير كافٍ.

\section{الحدود المفتوحة}

النتيجة المثبتة هنا حد علوي كلاسيكي فقط. تبقى خارج نطاق الفصل:
\begin{enumerate}
  \item الصيغة التقاربية الدقيقة مع حد رئيسي ثابت؛
  \item مبرهنة مونتغمري--هولي وصيغها الموحدة؛
  \item مجالات أصغر لـ\(Q\) لا يغطيها فصل الموصلات الحالي؛
  \item النسخ في الفترات القصيرة؛
  \item التعميمات إلى متتاليات أو معاملات أخرى؛
  \item تحليل PVG لمتوسطات التباين على ألياف البواقي والموصلات.
\end{enumerate}

\section{خلاصة الفصل}

يتكون البرهان من خمس حركات مترابطة:
\begin{enumerate}
  \item تحويل تباين الفئات إلى متوسط تربيعي على الشخصيات؛
  \item رد الشخصيات إلى موصلاتها مع وزن \(W_Q(r)\) وتصحيحات محلية صغيرة؛
  \item تطبيق الغربال الكبير بوزن الموصل؛
  \item فصل الموصلات الصغيرة والكبيرة عند \(Y=(x/Q)\log(2Q)\)؛
  \item إعادة التجميع وامتصاص كلفة الاستحثاث.
\end{enumerate}
الناتج هو
\[
V_\psi(x,Q)\ll_A xQ\log x
\]
في المجال \(x(\log x)^{-A}\le Q\le x\)، مع ثابت غير فعال بسبب
سيغل--فالفش وحدها. وقد بقيت الصياغات الأقوى والمسار العام لباربان مؤجلة
بوضوح بدل خلطها بالحد المثبت.
